\documentclass[10pt,a4paper]{article}
\usepackage{amsmath,amssymb,amsfonts}
\usepackage{geometry}
\geometry{
    top=20mm,
    bottom=20mm,
    left=16mm,
    right=16mm,
    headheight=14pt,
    headsep=15pt,
    footskip=15pt
}
\usepackage{iftex}
\ifXeTeX
  \usepackage{fontspec}
\else
  \usepackage[utf8]{inputenc}
  \usepackage[T5,T1]{fontenc}
\fi
\usepackage[vietnamese]{babel}
\usepackage{xcolor}
\usepackage{tcolorbox}
\tcbuselibrary{skins,breakable}
\usepackage{fancyhdr}
\usepackage{newtxtext,newtxmath}

% Định nghĩa hệ màu sắc nhận diện của sách Stewart Calculus
\definecolor{stewartcyan}{RGB}{0, 130, 185}
\definecolor{stewartred}{RGB}{204, 34, 41}
\definecolor{stewartgreen}{RGB}{0, 133, 66}
\definecolor{stewartgray}{RGB}{110, 110, 110}

% Khung định nghĩa chuẩn Stewart
\newtcolorbox{definitionbox}{
    enhanced,
    colback=white,
    colframe=stewartred,
    boxrule=1.1pt,
    arc=0pt,
    left=8pt,
    right=8pt,
    top=6pt,
    bottom=6pt,
    boxsep=0pt
}

% Khung tiểu sử Leibniz bên lề
\newtcolorbox{leibnizbox}{
    enhanced,
    colback=white,
    colframe=stewartgreen,
    boxrule=0.8pt,
    arc=0pt,
    left=5pt,
    right=5pt,
    top=5pt,
    bottom=5pt,
    boxsep=0pt
}

% Thiết lập tiêu đề đầu trang và chân trang
\pagestyle{fancy}
\fancyhf{}
\renewcommand{\headrulewidth}{0pt}
\fancyhead[L]{\fontsize{9pt}{11pt}\selectfont\textbf{156}\qquad\textbf{\textsf{CHƯƠNG 2}}\quad\textsf{Giới hạn và Đạo hàm}}
\fancyfoot[C]{\fontsize{7.5pt}{9pt}\selectfont\color{stewartgray}Bản dịch số hóa phục vụ học tập và nghiên cứu. Nguyên tác: James Stewart, \textit{Calculus: Early Transcendentals}, 9th Edition.}

\setlength{\parindent}{0pt}
\setlength{\parskip}{0pt}

\begin{document}

\noindent
\begin{minipage}[t]{0.26\textwidth}
    \vspace{0pt}
    % Dóng hàng chính xác với dòng tính toán thứ hai trong lời giải Ví dụ 4
    \vspace*{34pt}
    \begin{center}
    {\color{stewartcyan}\fontsize{9pt}{11pt}\selectfont
    $\displaystyle \frac{\frac{a}{b} - \frac{c}{d}}{e} = \frac{ad - bc}{bd} \cdot \frac{1}{e}$
    }
    \end{center}

    % Dóng hàng ngang dòng thứ năm của lời giải Ví dụ 4
    \vspace{24pt}
    \begin{leibnizbox}
    {\noindent\bfseries\textsf{\color{stewartgreen}\normalsize Leibniz}}\par\vspace{4pt}
    {\fontsize{7.5pt}{9.7pt}\selectfont
    \noindent Gottfried Wilhelm Leibniz sinh năm 1646 tại Leipzig và theo học luật, thần học, triết học và toán học tại trường đại học ở đó, tốt nghiệp cử nhân năm 17 tuổi. Sau khi nhận bằng tiến sĩ luật ở tuổi 20, Leibniz bước vào ngành ngoại giao và dành phần lớn cuộc đời đi đến các thủ đô của châu Âu để thực hiện các sứ mệnh chính trị. Đặc biệt, ông đã nỗ lực ngăn chặn mối đe dọa quân sự của Pháp đối với Đức và cố gắng hòa giải giữa Giáo hội Công giáo và Giáo hội Kháng cách.

    \par\vspace{4pt}
    \noindent\hspace{0.8em}Việc nghiên cứu toán học nghiêm túc của ông chỉ thực sự bắt đầu vào năm 1672 khi ông đang làm nhiệm vụ ngoại giao tại Paris. Tại đây, ông đã chế tạo một máy tính cơ học và gặp gỡ các nhà khoa học, như Huygens, người đã hướng sự chú ý của ông đến những phát triển mới nhất trong toán học và khoa học. Leibniz tìm cách xây dựng một hệ thống logic ký hiệu và ký pháp nhằm đơn giản hóa việc lập luận logic. Đặc biệt, phiên bản giải tích mà ông công bố vào năm 1684 đã thiết lập ký hiệu và các quy tắc tính đạo hàm mà chúng ta vẫn sử dụng ngày nay.

    \par\vspace{4pt}
    \noindent\hspace{0.8em}Thật không may, một cuộc tranh chấp quyền ưu tiên gay gắt đã nổ ra vào những năm 1690 giữa những người ủng hộ Newton và những người ủng hộ Leibniz về việc ai là người đầu tiên phát minh ra giải tích. Leibniz thậm chí còn bị các thành viên của Hội Hoàng gia Anh cáo buộc đạo văn. Sự thật là mỗi người đều đã phát minh ra giải tích một cách độc lập. Newton đã tìm ra phiên bản giải tích của mình trước, nhưng vì e ngại tranh cãi nên ông đã không công bố ngay. Do đó, công trình giải tích năm 1684 của Leibniz là công trình đầu tiên được xuất bản.
    }
    \end{leibnizbox}
\end{minipage}%
\hfill
\begin{minipage}[t]{0.71\textwidth}
    \vspace{0pt}
    \noindent\textbf{\textsf{\color{stewartcyan}VÍ DỤ 4}}\hspace{1em}Tìm $f'$ nếu $f(x) = \dfrac{1 - x}{2 + x}$.

    \vspace{0.6em}
    \noindent\textbf{\textsf{\color{stewartcyan}LỜI GIẢI}}

    \vspace{0.3em}
    \noindent\hspace{1.5em}$\begin{aligned}
    f'(x) &= \lim_{h \to 0} \frac{f(x + h) - f(x)}{h} \\[4pt]
          &= \lim_{h \to 0} \frac{\dfrac{1 - (x + h)}{2 + (x + h)} - \dfrac{1 - x}{2 + x}}{h} \\[8pt]
          &= \lim_{h \to 0} \frac{(1 - x - h)(2 + x) - (1 - x)(2 + x + h)}{h(2 + x + h)(2 + x)} \\[8pt]
          &= \lim_{h \to 0} \frac{(2 - x - 2h - x^2 - xh) - (2 - x + h - x^2 - xh)}{h(2 + x + h)(2 + x)} \\[8pt]
          &= \lim_{h \to 0} \frac{-3h}{h(2 + x + h)(2 + x)} \\[8pt]
          &= \lim_{h \to 0} \frac{-3}{(2 + x + h)(2 + x)} = -\frac{3}{(2 + x)^2} \qquad\quad {\color{stewartcyan}\blacksquare}
    \end{aligned}$

    \vspace{1.2em}
    \noindent{\color{stewartcyan}\rule{1.1ex}{1.1ex}}\hspace{0.5em}\textbf{\textsf{\large Các ký hiệu khác}}

    \vspace{0.4em}
    \noindent Nếu chúng ta sử dụng ký hiệu truyền thống $y = f(x)$ để chỉ ra rằng biến độc lập là $x$ và biến phụ thuộc là $y$, thì một số ký hiệu thay thế thông dụng cho đạo hàm như sau:
    \[
    f'(x) = y' = \frac{dy}{dx} = \frac{df}{dx} = \frac{d}{dx}f(x) = Df(x) = D_x f(x)
    \]

    \noindent Các ký hiệu $D$ và $d/dx$ được gọi là các \textbf{toán tử đạo hàm} vì chúng chỉ ra phép toán \textbf{lấy đạo hàm}, tức là quá trình tính toán một đạo hàm.

    \vspace{0.4em}
    \noindent Ký hiệu $dy/dx$, do Leibniz đưa ra, (vào thời điểm hiện tại) không nên được coi là một tỉ số; nó đơn giản chỉ là một cách viết đồng nghĩa với $f'(x)$. Tuy nhiên, đây là một ký hiệu rất hữu ích và mang tính gợi mở, đặc biệt khi được sử dụng kết hợp với ký hiệu số gia. Dựa vào Phương trình 2.7.6, chúng ta có thể viết lại định nghĩa đạo hàm theo ký hiệu Leibniz dưới dạng
    \[
    \frac{dy}{dx} = \lim_{\Delta x \to 0} \frac{\Delta y}{\Delta x}
    \]

    \noindent Nếu chúng ta muốn biểu thị giá trị của đạo hàm $dy/dx$ theo ký hiệu Leibniz tại một số cụ thể $a$, ta dùng ký hiệu
    \[
    \left.\frac{dy}{dx}\right|_{x=a} \qquad \text{hoặc} \qquad \left.\frac{dy}{dx}\right]_{x=a}
    \]
    ký hiệu này đồng nghĩa với $f'(a)$. Dấu gạch thẳng đứng có nghĩa là ``tính giá trị tại''.

    \vspace{0.8em}
    \begin{definitionbox}
    \noindent\setlength{\fboxrule}{0.9pt}\setlength{\fboxsep}{2.5pt}\fcolorbox{stewartred}{white}{\bfseries\color{stewartred}\strut 3}\hspace{0.6em}{\bfseries\color{stewartred}\textsf{Định nghĩa}}\hspace{0.8em}Một hàm số $f$ được gọi là \textbf{khả vi tại $a$} nếu $f'(a)$ tồn tại. Hàm số được gọi là \textbf{khả vi trên một khoảng mở} $(a, b)$ [hoặc $(a, \infty)$ hoặc $(-\infty, a)$ hoặc $(-\infty, \infty)$] nếu nó khả vi tại mọi số trong khoảng đó.
    \end{definitionbox}
\end{minipage}

\end{document}